\section*{第 \ref{ch:b2:realfun} 章——实变函数}

\begin{solution}{exo:b2:realfun:1}
$\lfloor x\rfloor$：在每个整数处跳跃，跳跃量为 $1$。$x -
\lfloor x\rfloor$：在整数处跳跃，跳跃量为 $-1$（左极限为 $1$，
值为 $0$）。$\lfloor x\rfloor + \sqrt{x - \lfloor x\rfloor}$：在整数
$n$ 处，左极限为 $(n - 1) + 1 = n$，取值为 $n$：
处处\emph{\omterm{def:b2:metric:continuity}{连续}}（平方根弥补了跳跃），
但在整数处不可微。有理数跳跃函数：
将构造限制在二进有理数的一个枚举上，
其跳跃量为 $2^{-n}$，且恰在第 $n$ 个二进有理数处
发生跳跃，其余各处\omterm{def:b2:metric:continuity}{连续}。
\end{solution}

\begin{solution}{exo:b2:realfun:2}
假设递增函数 $f$ 在内点 $a$ 处有间断，则
$f(a^-) < f(a^+)$（见\cref{thm:b2:realfun:monotone}），且
$I$ 的像避开非空开区间
$\intoo{f(a^-)}{f(a^+)}$，至多可能除去单点 $f(a)$：
包含 $a$ 为内点的任一子区间的像都
不是区间（在 $f(a)$ 的至少一侧有空隙）。
这与介值性质相矛盾。端点处的间断
同理可借助单侧空隙予以排除。
\end{solution}

\begin{solution}{exo:b2:realfun:3}
$x\ln x$：二阶导数 $\frac1x > 0$，故凸。$\ln(1 +
\eu^x)$：导数 $\frac{\eu^x}{1 + \eu^x} = 1 - \frac{1}{1 +
\eu^x}$，递增，故凸。$\sqrt{1 + x^2}$：二阶导数
$(1 + x^2)^{-3/2} > 0$，故凸。$x^3$：在 $\R$ 上不凸（$f'' =
6x$ 变号）；仅在 $\R_+$ 上凸。
\end{solution}

\begin{solution}{exo:b2:realfun:4}
假设 $f(a) \neq f(b)$，不妨设 $f(b) > f(a)$，其中 $a < b$（另一情形
$f(b) < f(a)$ 同理，向左考虑）。对 $x > b$，由斜率不等式
（\cref{lem:b2:realfun:slopes}）在 $a < b < x$ 上给出
\[
\frac{f(x) - f(a)}{x - a} \geq \frac{f(b) - f(a)}{b - a} = m > 0
\quad\Longrightarrow\quad
f(x) \geq f(a) + m(x - a) \xrightarrow[x\to+\infty]{} +\infty,
\]
与有上界性相矛盾。因此 $f$ 为常数。

渐近线：若 $f(x) - (\alpha x + \beta) \to 0$ 在 $+\infty$ 处成立，且
$f(x) - (\alpha' x + \beta') \to 0$ 在 $-\infty$ 处成立，记凸函数
$g(x) = f(x) - (\alpha x + \beta)$ 在
$+\infty$ 附近有上界；凸性加上 $-\infty$ 处的渐近线（这迫使
$\alpha' \leq \alpha$，继而通过比较斜率得 $\alpha' = \alpha$
（在 $\mp\infty$ 处比较，凸函数的斜率递增），使得 $g$
在整个 $\R$ 上有上界，从而在极限意义下为常数 $= 0$，故 $f$
为仿射映射。
\end{solution}

\begin{solution}{exo:b2:realfun:5}
$f'$ 单调，由\cref{thm:b2:realfun:monotone}可知其
唯一可能的间断点为跳跃型，且单侧极限处处存在。
由\cref{cor:b2:realfun:nojumps}，导数没有
跳跃型间断点。因此 $f'$ 没有任何间断点，
即\omterm{def:b2:metric:continuity}{连续}。
\end{solution}

\begin{solution}{exo:b2:realfun:6}
取对数：
\[
\frac1p \ln\Bigl(\sum_i \lambda_i x_i^p\Bigr)
= \frac1p \ln\Bigl(1 + p\sum_i \lambda_i \ln x_i +
O(p^2)\Bigr)
= \sum_i \lambda_i \ln x_i + O(p)
\xrightarrow[p \to 0^+]{} \sum_i \lambda_i \ln x_i ,
\]
利用 $\sum\lambda_i = 1$ 及 $\ln(1 + u) = u + O(u^2)$。
取指数即得几何平均。现对每个 $p \in
\intoo{0}{1}$，由幂平均不等式
（\cref{ex:b2:realfun:powermeans}，指数 $p < 1$）给出
\[
\Bigl(\sum_i \lambda_i x_i^p\Bigr)^{1/p} \leq \sum_i\lambda_i x_i ;
\]
令 $p \to 0^+$，在左端得 $\prod_i x_i^{\lambda_i}
\leq \sum_i \lambda_i x_i$：即加权算术--几何平均不等式。
\end{solution}

\begin{solution}{exo:b2:realfun:7}
取 $\varphi(t) = t\ln t$（凸：$\varphi'' = \frac1t > 0$），并以
$q_i$ 为权，取点 $t_i = \frac{p_i}{q_i}$：
\[
\sum_i p_i \ln\frac{p_i}{q_i}
= \sum_i q_i\, \varphi\Bigl(\frac{p_i}{q_i}\Bigr)
\;\geq\; \varphi\Bigl(\sum_i q_i \frac{p_i}{q_i}\Bigr)
= \varphi(1) = 0 ,
\]
由琴生不等式（\cref{thm:b2:realfun:convexreg}~(3)）可得。琴生不等式取等号时，
对\emph{严格}凸函数，取等号迫使所有点
$t_i$ 重合：$\frac{p_i}{q_i}$ 为常数，求和知该常数为
$1$，故 $p = q$。（此量——即库尔贝克--莱布勒
散度——在\cref{ch:b2:randomvar}中再次出现。）
\end{solution}

\begin{solution}{exo:b2:realfun:8}
\emph{二进权重。}对 $m$ 作归纳：$m = 1$ 的情形即为
假设。对权重 $\lambda = \frac{k}{2^{m+1}}$（$k$ 为奇数），
记 $\lambda = \frac12(\lambda_1 + \lambda_2)$，其中 $\lambda_j =
\frac{k \mp 1}{2^{m+1}}$，两者化简后分母均为 $2^m$；
于是
\[
f\bigl(\lambda x + (1{-}\lambda)y\bigr)
= f\Bigl(\tfrac{u + v}{2}\Bigr)
\leq \frac{f(u) + f(v)}{2}
\leq \lambda f(x) + (1 - \lambda) f(y),
\]
其中 $u = \lambda_1 x + (1 - \lambda_1)y$，且 $v = \lambda_2 x +
(1-\lambda_2)y$，先用中点凸性，再对
$u, v$ 用归纳假设。

\emph{取极限。}对任意 $\lambda \in
\intcc{0}{1}$，取二进有理数列 $\lambda_n \to \lambda$，
由 $f$ 与仿射映射的\omterm{def:b2:metric:continuity}{连续}性，上述不等式
$f(\lambda_n x + (1-\lambda_n)y) \leq \lambda_n f(x) +
(1-\lambda_n)f(y)$ 在极限下成立，故 $f$ 为凸函数。
\end{solution}

\begin{solution}{exo:b2:realfun:9}
对 $0 < x < y$，由斜率不等式
（见\cref{lem:b2:realfun:slopes}）在点 $0 < x < y$ 处给出
\[
\frac{f(x) - f(0)}{x} \leq \frac{f(y) - f(0)}{y},
\qquad\text{即}\qquad
\frac{f(x)}{x} \leq \frac{f(y)}{y} +
f(0)\Bigl(\frac1x - \frac1y\Bigr).
\]
由于 $f(0) \leq 0$ 且 $\frac1x - \frac1y > 0$，末项
$\leq 0$，故 $\frac{f(x)}{x} \leq \frac{f(y)}{y}$。于是 $x \mapsto
\frac{f(x)}x$ 递增。

次可加性：设 $f(0) = 0$，对 $x, y > 0$（有零变量的情形
平凡），
\[
f(x) = x\,\frac{f(x)}{x} \leq x\,\frac{f(x+y)}{x+y},
\qquad
f(y) \leq y\,\frac{f(x+y)}{x+y},
\]
由刚证明的单调性，相加得 $f(x) + f(y) \leq
f(x+y)$。
\end{solution}

\begin{solution}{exo:b2:realfun:10}
对 $x, y \in I$ 及 $\lambda \in \intcc01$，先用 $f$ 的凸性，
再用 $g$ 的单调性，最后用 $g$ 的凸性：
\[
g\bigl(f(\lambda x + (1{-}\lambda)y)\bigr)
\leq g\bigl(\lambda f(x) + (1{-}\lambda)f(y)\bigr)
\leq \lambda\,g(f(x)) + (1{-}\lambda)\,g(f(y)).
\]
反例（去掉单调性）：$g(t) = -t$ 为凸函数
（仿射的）但递减，$f(x) = x^2$ 为凸函数，而 $g \circ f =
-x^2$ 为严格凹函数。
\end{solution}

\begin{solution}{exo:b2:realfun:11}
\emph{左侧不等式：}设 $m = \frac{a+b}2$，在
$m$ 处取一条支撑线（见\cref{thm:b2:realfun:convexreg}~(2)）：$f(t) \geq
f(m) + \mu(t - m)$，对所有 $t \in \intcc ab$ 成立。在
$\intcc{a}{b}$ 上积分：线性项的积分为 $\mu\int_a^b(t -
m)\dd t = 0$（关于 $m$ 对称），故 $\int_a^b f \geq (b -
a)f(m)$。

\emph{右侧不等式：}在 $\intcc ab$ 上，由凸性，$f$ 以
其弦为上界：$f(t) \leq f(a) + \frac{f(b) - f(a)}{b - a}(t - a)$。
积分得 $\int_a^b f \leq (b-a)f(a) + \frac{f(b) - f(a)}{b -
a}\cdot\frac{(b-a)^2}2 = (b - a)\,\frac{f(a) + f(b)}2$。两边除以
$b - a$。
\end{solution}

\begin{solution}{exo:b2:realfun:12}
设 $a - \delta < a \leq x < y \leq b < b + \delta$，均属于 $I$。
两次应用斜率不等式
（见\cref{lem:b2:realfun:slopes}），先对 $a - \delta < a \leq x <
y$，再对 $x < y \leq b < b + \delta$：
\[
\frac{f(a) - f(a - \delta)}{\delta}
\leq \frac{f(y) - f(x)}{y - x}
\leq \frac{f(b + \delta) - f(b)}{\delta}
\]
（两端点同时右移时，弦的斜率增大）。由此
$\intcc ab$ 内任一弦的斜率都夹在两个固定
数之间，于是
\[
\abs{f(y) - f(x)} \leq K\,\abs{y - x},
\qquad
K = \max\Bigl(\Bigl|\frac{f(a) - f(a-\delta)}{\delta}\Bigr|,
\Bigl|\frac{f(b+\delta) - f(b)}{\delta}\Bigr|\Bigr):
\]
$f$ 在 $\intcc ab$ 上满足\omterm{def:b2:metric:continuity}{利普希茨}条件。开区间 $I$ 的每一点
周围都有这样一段带余量的线段，因而局部\omterm{def:b2:metric:continuity}{利普希茨}，
从而在 $I$ 上\omterm{def:b2:metric:continuity}{连续}。
\end{solution}

\begin{solution}{pb:b2:realfun:1}
\textbf{1.} \emph{凸 $\Rightarrow$ $f'$ 递增：}对 $a <
b$，由斜率不等式，对充分小的 $h > 0$，有 $\frac{f(a+h)
- f(a)}h \leq \frac{f(b) - f(a)}{b-a} \leq \frac{f(b) -
f(b-h)}{h}$，令 $h \to 0$：$f'(a) \leq \frac{f(b) -
f(a)}{b - a} \leq f'(b)$。\emph{反之}，若 $f'$ 递增，
且 $x < y < z$，由中值定理，存在 $c_1 \in
\intoo{x}{y}$，$c_2 \in \intoo yz$，使得
\[
\frac{f(y) - f(x)}{y - x} = f'(c_1) \leq f'(c_2)
= \frac{f(z) - f(y)}{z - y},
\]
由此三点斜率不等式，取 $y = \lambda
x + (1 - \lambda)z$，整理即得凸性不等式。
对 $C^2$ 函数，$f'' \geq 0$ 当且仅当 $f'$ 递增。

\textbf{2.} 若 $f'' > 0$，则 $f'$ 严格递增，
由上面的中值计算，弦斜率之间得到一个\emph{严格}不等式，
即严格凸性。\emph{严格
支撑：}在内点 $a$ 处，其支撑斜率为 $m$，若 $f(x_0) =
f(a) + m(x_0 - a)$，对某个 $x_0 \neq a$ 成立，则在
从 $a$ 到 $x_0$ 的线段上，支撑线与弦重合，并且
由中点处的严格凸性得 $f\bigl(\frac{a +
x_0}2\bigr) < f(a) + m\,\frac{x_0 - a}2$，与
支撑不等式矛盾。故 $f(x) > f(a) + m(x - a)$，对所有 $x \neq
a$。\emph{严格詹森不等式：}取 $a = \sum\lambda_ix_i$，对
支撑不等式取平均，得 $\sum\lambda_if(x_i) \geq f(a)$，
等号成立当且仅当每一项都取等号，即当且仅当每个 $x_i
= a$。

\textbf{3.} $(-\ln)'' = \frac1{t^2} > 0$; $(t^r)'' = r(r -
1)t^{r-2}$，当 $r > 1$ 时为正，当 $0 < r < 1$ 时为负；
$\exp'' = \exp > 0$。由第 2 题，以上均为严格不等式。

\textbf{4.} $ab = 0$ 的情形显然。对 $a, b > 0$，
由 $\ln$ 在点 $a^p, b^q$ 处的凹性，权重分别取为
$\frac1p, \frac1q$：
\[
\ln\Bigl(\frac{a^p}p + \frac{b^q}q\Bigr) \geq
\frac1p\ln(a^p) + \frac1q\ln(b^q) = \ln(ab),
\]
又 $\ln$ 单调递增，故 $ab \leq \frac{a^p}p + \frac{b^q}q$。
等号成立当且仅当两点重合（由严格凹性）：$a^p =
b^q$。

\textbf{5.} 对 $\ln$ 取权重 $\lambda_i$，应用凹性：
$\ln\bigl(\sum\lambda_ix_i\bigr) \geq \sum\lambda_i\ln x_i =
\ln\prod x_i^{\lambda_i}$；再取指数。等号成立当且仅当所有 $x_i$ 
相等（由第 2 题）。\cref{exo:b2:realfun:6} 的思路
通过幂平均的极限得到了同一不等式；此处
只需应用一次琴生不等式——
这套工具本身即含冗余。

\textbf{6.} 若 $a = 0$ 或 $b = 0$，不等式显然成立。
归一化：将 $a$ 换为 $a/\norm a_p$，$b$ 换为 $b/\norm
b_q$，可设 $\norm a_p = \norm b_q = 1$，并需证明
$\sum\abs{a_ib_i} \leq 1$。逐项应用杨不等式：
\[
\sum_i\abs{a_i}\abs{b_i} \leq
\sum_i\Bigl(\frac{\abs{a_i}^p}{p} +
\frac{\abs{b_i}^q}{q}\Bigr) = \frac1p + \frac1q = 1 .
\]
等号成立当且仅当每个杨不等式都取等号：$\abs{a_i}^p =
\abs{b_i}^q$ 对所有 $i$ 成立——还原归一化后，
即 $(\abs{a_i}^p)$ 与 $(\abs{b_i}^q)$ 成比例。

\textbf{7.} $p = q = 2$ 即柯西--施瓦茨不等式，
等号条件相同（成比例）。端点情形：$\sum\abs{a_ib_i}
\leq \bigl(\max_i\abs{b_i}\bigr)\sum_i\abs{a_i} =
\norm a_1\norm b_\infty$，逐项即得。

\textbf{8.} 当 $p = 1$ 时，逐项即为三角不等式。
当 $p > 1$ 时，记 $q$ 为其共轭指数：
\[
\norm{a+b}_p^p = \sum_i\abs{a_i + b_i}^p
\leq \sum_i\abs{a_i+b_i}^{p-1}\abs{a_i} +
\sum_i\abs{a_i+b_i}^{p-1}\abs{b_i},
\]
对每个和式应用赫尔德不等式，注意 $(p - 1)q = p$：
\[
\sum_i\abs{a_i+b_i}^{p-1}\abs{a_i} \leq
\Bigl(\sum_i\abs{a_i+b_i}^{p}\Bigr)^{1/q}\norm a_p
= \norm{a + b}_p^{p/q}\,\norm a_p ,
\]
对 $b$ 同样处理。因此 $\norm{a+b}_p^p \leq \norm{a +
b}_p^{p/q}\bigl(\norm a_p + \norm b_p\bigr)$；若 $a + b \neq 0$，
两边除以 $\norm{a+b}_p^{p/q}$，并利用 $p - \frac pq = 1$。
由齐次性与分离性（显然），$\norm\cdot_p$ 是
$\R^n$ 上的\omterm{def:b2:nvs:norm}{范数}。

\textbf{9.} 设 $f, g$ 为 $\intcc ab$ 上的\omterm{def:b2:metric:continuity}{连续}函数，赫尔德不等式为
\[
\int_a^b\abs{fg} \leq \Bigl(\int_a^b\abs
f^p\Bigr)^{1/p}\Bigl(\int_a^b\abs g^q\Bigr)^{1/q}
\]
由相同的归一化与逐点杨不等式，积分后得到；
同样地，闵可夫斯基不等式 $\norm{f + g}_p \leq \norm f_p + \norm g_p$ 由
相同的拆分，并对每一部分应用赫尔德不等式。\omterm{def:b2:nvs:norm}{范数}的分离性
依赖严格正性：\omterm{def:b2:metric:continuity}{连续}函数 $\abs f^p$ 的积分为零时
必恒为零（见本科一年级分册）。

\textbf{10.} \emph{单调性：}可设 $\norm a_p = 1$；
则各 $\abs{a_i} \leq 1$，故 $\abs{a_i}^q \leq \abs{a_i}^p$
且 $\norm a_q^q \leq 1$，即 $\norm a_q \leq 1 = \norm a_p$。
等号成立要求 $\abs{a_i}^q = \abs{a_i}^p$ 对每个 $i$ 成立，
即每个 $\abs{a_i} \in \{0, 1\}$；再由 $\sum\abs{a_i}^p = 1$，
恰有一个分量的模为 $1$：等号成立当且仅当
$a$ 至多有一个非零分量。\emph{极限：}$\norm a_\infty \leq
\norm a_p \leq n^{1/p}\norm a_\infty$，且 $n^{1/p} \to 1$。
\emph{反向比较：}赫尔德不等式，指数取为 $\frac
qp$ 及其共轭 $\frac{q}{q-p}$，作用于
$\abs{a_i}^p\cdot 1$：
\[
\norm a_p^p = \sum_i\abs{a_i}^p\cdot 1 \leq
\Bigl(\sum_i\abs{a_i}^{q}\Bigr)^{p/q}\,n^{1 - p/q}
= \norm a_q^{p}\; n^{1-p/q},
\]
由此 $\norm a_p \leq n^{\frac1p - \frac1q}\norm a_q$，
等号成立当且仅当所有 $\abs{a_i}$ 相等（即赫尔德不等式
对常向量的等号情形）。

\textbf{11.} 写 $\abs{a_i}^r = \abs{a_i}^{\theta
r}\,\abs{a_i}^{(1-\theta)r}$，并应用赫尔德不等式，取
共轭指数 $\frac{p}{\theta r}$ 与
$\frac{q}{(1-\theta)r}$（二者确为共轭，因为 $\frac{\theta
r}p + \frac{(1-\theta)r}q = 1$）：
\[
\norm a_r^r = \sum_i \abs{a_i}^{\theta r}\abs{a_i}^{(1-\theta)r}
\leq \Bigl(\sum_i\abs{a_i}^{p}\Bigr)^{\theta r/p}
\Bigl(\sum_i\abs{a_i}^{q}\Bigr)^{(1-\theta)r/q}
= \norm a_p^{\theta r}\,\norm a_q^{(1-\theta)r} .
\]
两边开 $r$ 次方：$p$ \omterm{def:b2:nvs:norm}{范数}关于 $\frac1p$ 为对数凸。

\textbf{12.} \emph{两者皆为负：}若 $p < q < 0$，则 $0 < -q <
-p$，对正指数有 $M_{-q}(y) \leq M_{-p}(y)$，
（课程中的情形，见\cref{ex:b2:realfun:powermeans}）应用于 $y =
(1/x_i)$；对恒等式 $M_p(x) = M_{-p}(1/x)^{-1}$ 取逆，
不等号随之反向，即得 $M_p(x) \leq M_q(x)$。
\emph{桥接：}当 $q > 0$ 时，$\ln$ 的凹性给出 $\ln M_q =
\frac1q\ln\bigl(\sum\lambda_ix_i^q\bigr) \geq
\frac1q\sum\lambda_i\ln x_i^q = \ln M_0$；当 $p < 0$ 时，同样
由凹性给出 $\ln\bigl(\sum\lambda_ix_i^p\bigr) \geq
p\sum\lambda_i\ln x_i$，除以 $p < 0$ 后不等号反向：$\ln M_p
\leq \ln M_0$。因此 $M_p \leq M_0 \leq M_q$，只要 $p < 0 <
q$ 即成立；连同两种同号情形，$M$ 在整个 
$\R^*$ 上（并经过 $0$）单调递增。

\textbf{13.} 设 $x_{\max} = \max x_i$，该最大值在 $i^*$ 处取得。对
$p > 0$，有
\[
\lambda_{i^*}^{1/p}\,x_{\max} \leq M_p \leq x_{\max},
\]
且 $\lambda_{i^*}^{1/p} \to 1$，于是 $M_p \to x_{\max}$。当 $p \to
-\infty$ 时：$M_p(x) = M_{-p}(1/x)^{-1} \to
\bigl(\max_i\frac1{x_i}\bigr)^{-1} = \min_ix_i$。

\textbf{14.} 取 $\lambda_i = \frac1n$，不等式链 $M_{-\infty}
\leq M_{-1} \leq M_0 \leq M_1 \leq M_2 \leq M_{+\infty}$ 可写为
\[
\min \leq \frac{n}{\sum\frac1{a_i}} \leq \Bigl(\prod
a_i\Bigr)^{1/n} \leq \frac{\sum a_i}{n} \leq
\sqrt{\frac{\sum a_i^2}{n}} \leq \max .
\]
算术–调和平均不等式不等式（$M_{-1} \leq M_1$）直接整理为
$\bigl(\sum a_i\bigr)\bigl(\sum\frac1{a_i}\bigr) \geq n^2$。

\textbf{15.} 等权重时，$M_p(x) =
\bigl(\frac1n\sum\abs{x_i}^p\bigr)^{1/p} = n^{-1/p}\norm x_p$。
随 $p$ 增大，$\norm x_p$ 减小（第 10 题），但
归一化因子 $n^{-1/p}$ 增长更快，于是乘积
增大（第 12 题）：均值取平均，\omterm{def:b2:nvs:norm}{范数}则是累积，而
因子 $n^{-1/p}$ 正是这两种记账约定之间的
换算比率。

\textbf{16.} 每一环节都是严格琴生不等式
（第 2 题）的实例，所用函数为第 3 题中的严格凸、凹函数，
即 $t^{q/p}$ 与 $\ln$，故任一环节取等号都迫使
所有 $x_i$ 相等；又 $\min = M_p$ 或 $M_p = \max$ 同样
迫使所有取值都等于该公共极值。
只要有两个 $x_i$ 不等，整条链便是严格的。

\textbf{17.} 将杨氏不等式（第 4 题）应用于数对
$\varepsilon^{1/p}a$ 与 $\varepsilon^{-1/p}b$：
\[
ab = (\varepsilon^{1/p}a)(\varepsilon^{-1/p}b)
\leq \varepsilon\,\frac{a^p}p +
\varepsilon^{-q/p}\,\frac{b^q}q .
\]
取 $p = q = 2$，将 $\varepsilon$ 换为 $2\varepsilon$：
$ab \leq \varepsilon a^2 + \frac{b^2}{4\varepsilon}$——这就是
吸收不等式：一个乘积换成一个平方的小倍数，
再加上另一个平方的大倍数。

\textbf{18.} 裂项相消：
\[
\prod_{k=1}^{n}c_k = \frac{\prod_{k=1}^n(k+1)^k}
{\prod_{k=1}^{n}k^{k-1}}
= \frac{2^1\,3^2\cdots(n+1)^n}{1^0\,2^1\cdots n^{n-1}}
= (n+1)^n,
\]
分子中每个因子 $(k+1)^k$ 都与分母中的
下一项相消。对 $n$ 个数 $c_ka_k$ 应用算术–几何平均不等式：
\[
(a_1\cdots a_n)^{1/n} =
\frac{\bigl(\prod_k c_ka_k\bigr)^{1/n}}{(n+1)}
\leq \frac{1}{n+1}\cdot\frac1n\sum_{k=1}^{n}c_ka_k .
\]

\textbf{19.} 对 $n$ 求和，并交换两重求和的次序
（各项均为正，故可用\cref{thm:b2:series:fubini}）：
\[
\sum_{n\geq1}(a_1\cdots a_n)^{1/n}
\leq \sum_{n\geq1}\frac{1}{n(n+1)}\sum_{k=1}^{n}c_ka_k
= \sum_{k\geq1}c_ka_k\sum_{n\geq k}\frac1{n(n+1)}
= \sum_{k\geq1}\frac{c_ka_k}{k},
\]
利用裂项 $\sum_{n\geq k}\bigl(\frac1n -
\frac1{n+1}\bigr) = \frac1k$。最后 $\frac{c_k}k =
\frac{(k+1)^k}{k^k} = \bigl(1 + \frac1k\bigr)^k < \eu$
（单调递增数列，极限为 $\eu$，见本科一年级分册）：
\[
\sum_{n\geq1}(a_1\cdots a_n)^{1/n} \leq
\eu\sum_{k\geq1}a_k :
\]
卡莱曼不等式。（常数 $\eu$ 最优，但
在此不予证明。）

\textbf{20.} 将柯西--施瓦茨不等式（第 9 题，$p = q = 2$）
应用于 $\sqrt f$ 与 $\frac1{\sqrt f}$：
\[
1 = \Bigl(\int_0^1\sqrt f\cdot\frac{1}{\sqrt f}\Bigr)^{2}
\leq \Bigl(\int_0^1 f\Bigr)\Bigl(\int_0^1\frac1f\Bigr).
\]
等号成立当且仅当 $\sqrt f$ 与 $\frac1{\sqrt f}$ 成比例，
即 $f^2$ 为常数，即 $f$ 为常数（$f > 0$ \omterm{def:b2:metric:continuity}{连续}）。

\textbf{21.} 设 $1 < p < \infty$，$\norm a_p = \norm b_p = 1$，
$a \neq b$，并假设 $\bigl\Vert\frac{a+b}2\bigr\Vert_p = 1$，
即对 $a, b$ 闵可夫斯基不等式取等号。沿第 8 题的
证明，则等号迫使两次赫尔德不等式的应用均取等号，
以及逐项三角不等式中的取等：$(\abs{a_i}^p)$ 与
$(\abs{b_i}^p)$ 都与 $(\abs{a_i + b_i}^p)$ 成比例，并且
$a_i, b_i$ 同号——因此 $b = ta$，其中 $t \geq
0$，且由 $\norm b_p = \norm a_p$ 得 $t = 1$：$b = a$，
矛盾。因此 $p$ 球面不含任何两个不同
球面点中点，从而不含线段。当 $p = \infty$ 时，在 $\R^2$ 中：
所有 $(1, t)$，$\abs t \leq 1$，都在单位球面上——是一段平直的
边；而当 $p = 1$ 时，线段 $(t, 1 - t)$，$t \in \intcc01$，
则确含线段。

\textbf{22.} 当 $a = 0$ 时两边均为零。否则由赫尔德不等式，
每个 $\sum a_ib_i$ 都受 $\norm a_p\norm b_q \leq \norm
a_p$ 为界。取等情形：取
\[
b_i = \frac{\operatorname{sign}(a_i)\,\abs{a_i}^{p-1}}
{\norm a_p^{p/q}} :
\qquad
\norm b_q^q = \frac{\sum_i\abs{a_i}^{(p-1)q}}{\norm a_p^{p}}
= \frac{\norm a_p^p}{\norm a_p^p} = 1,
\quad
\sum_ia_ib_i = \frac{\norm a_p^p}{\norm a_p^{p/q}} =
\norm a_p ,
\]
利用 $(p-1)q = p$ 与 $p - \frac pq = 1$。于是上确界
是最大值，且等于 $\norm a_p$：每个 $p$ \omterm{def:b2:nvs:norm}{范数}都是
其共轭指数的对偶\omterm{def:b2:nvs:norm}{范数}——这正是 $L^p$--$L^q$ 对偶性的萌芽。

\textbf{23.} $\E[X^r] = \sum_i\lambda_ix_i^r$，于是
$\E[X^r]^{1/r} = M_r(x; \lambda)$，由第 12 题知其关于 $r$ 递增
（$r \to 0, \pm\infty$ 的情形由第
12--13 题给出）：这是李雅普诺夫矩不等式，纯粹是关于
加权幂平均的陈述。对真正的随机变量，它会
\cref{ch:b2:randomvar}中再次出现。

\textbf{24.} (i) 对幂平均 $M_1 \leq M_3$，取等权：
$\frac{a+b+c}3 \leq \bigl(\frac{a^3+b^3+c^3}3\bigr)^{1/3}$；
两边立方并乘以 $3$：$a^3 + b^3 + c^3 \geq
\frac{(a+b+c)^3}9$。(ii) 对常值
向量 $\sum_i\sqrt{x_i}\cdot1 \leq
\bigl(\sum_ix_i\bigr)^{1/2}n^{1/2}$；两边平方。

\textbf{25.} 由弦的定义，经一步
代数变形即得斜率引理；在一点被夹逼的斜率给出
单侧导数与支撑线，它们的加权平均便是
即为琴生不等式。对 $-\ln$ 应用后，琴生不等式化为杨不等式，
对归一化向量求和即为赫尔德不等式，经拆分
再重新吸收即为闵可夫斯基不等式——由此在
\cref{ch:b2:nvs}中诞生了 $p$ \omterm{def:b2:nvs:norm}{范数}，连同其对偶性（第 22 题）
与其几何（第 21 题）。将琴生不等式沿幂的尺度
应用，串起从 $\min$ 到 $\max$ 的所有平均
（第 12--14 题），在随机变量上解读即为
矩不等式（第 23 题）。算术–几何平均不等式
经一个伸缩技巧加权，得到卡勒曼界，其
不可约常数 $\eu$（第 18--19 题）。高峰：
赫尔德–闵可夫斯基，以及卡勒曼界。终点是 $L^p$
空间（本科三年级分册），其奠基公理恰好是
第 6 题与第 8 题，只需把求和换成积分。
\end{solution}
